% =====================================================================
% Apartado 4: CAP y PACELC (2 min) -- Responsable: Elian Mendoza
% Solo usa paquetes que ya carga presentacion.tex (booktabs, biblatex).
% =====================================================================
\section{CAP y PACELC}

% ---------------------------------------------------------------------
\begin{frame}{Teorema CAP}
  Un sistema distribuido replicado no puede garantizar simultáneamente
  las tres propiedades siguientes.

  \bigskip
  \centering
  \small
  \renewcommand{\arraystretch}{1.4}
  \begin{tabular}{p{0.26\textwidth} p{0.62\textwidth}}
    \toprule
    \textbf{Propiedad} & \textbf{Definición} \\
    \midrule
    Consistencia (C) & Toda lectura refleja la escritura más reciente \\
    Disponibilidad (A) & Toda solicitud recibe una respuesta \\
    Tolerancia a\newline particiones (P) & El sistema opera aunque se pierda la
      comunicación entre nodos \\
    \bottomrule
  \end{tabular}

  \vfill
  \raggedright
  \footnotesize Fuentes: \cite[p.~172]{ahmed2023consistency},
  \cite[p.~2]{lee2023consistency}.
\end{frame}

% ---------------------------------------------------------------------
\begin{frame}{Comportamiento ante una partición}
  Si una réplica no responde dentro del tiempo de espera, el diseño
  determina la respuesta del sistema.

  \bigskip
  \centering
  \small
  \renewcommand{\arraystretch}{1.4}
  \begin{tabular}{p{0.20\textwidth} p{0.33\textwidth} p{0.33\textwidth}}
    \toprule
    & \textbf{Sistema CP} & \textbf{Sistema AP} \\
    \midrule
    Ante la partición & Rechaza la solicitud & Responde con el dato
      disponible \\
    Garantiza & Consistencia & Disponibilidad \\
    Sacrifica & Disponibilidad & Consistencia (el dato puede estar
      desactualizado) \\
    \bottomrule
  \end{tabular}

  \vfill
  \raggedright
  \footnotesize Fuentes: \cite[p.~172]{ahmed2023consistency},
  \cite[p.~51:4]{almeida2024crdt}.
\end{frame}

% ---------------------------------------------------------------------
\begin{frame}{Teorema PACELC}
  PACELC extiende CAP al funcionamiento normal de la red.

  \bigskip
  \centering
  \small
  \renewcommand{\arraystretch}{1.4}
  \begin{tabular}{p{0.30\textwidth} p{0.56\textwidth}}
    \toprule
    \textbf{Situación} & \textbf{Decisión de diseño} \\
    \midrule
    Con partición (P) & Disponibilidad (A) o consistencia (C) \\
    Sin partición (E) & Latencia (L) o consistencia (C) \\
    \bottomrule
  \end{tabular}

  \bigskip
  \raggedright
  La consistencia fuerte requiere coordinar réplicas mediante protocolos
  de consenso; el tiempo de respuesta aumenta con el número de nodos y la
  distancia entre ellos.

  \vfill
  \footnotesize Fuentes: \cite[p.~172]{ahmed2023consistency},
  \cite[p.~51:4]{almeida2024crdt}.
\end{frame}

% ---------------------------------------------------------------------
\begin{frame}{Spotify ante una partición: el caso de Padlock}
  Padlock es el servicio central de Spotify que gestiona las claves de
  cifrado de los datos personales de sus usuarios.

  \bigskip
  \small
  \renewcommand{\arraystretch}{1.35}
  \begin{tabular}{p{0.24\textwidth} p{0.64\textwidth}}
    \toprule
    \textbf{Aspecto} & \textbf{Documentación de Spotify} \\
    \midrule
    Almacenamiento & Cassandra, elegida por su replicación global y alta
      disponibilidad \\
    Lecturas & Nivel LOCAL\_ONE; reintento con LOCAL\_QUORUM solo si no
      encuentra la clave \\
    Despliegue & Tres regiones con tres zonas cada una \\
    Preparación & Pruebas que simulan caídas de centros de datos y
      divisiones de red \\
    \bottomrule
  \end{tabular}

  \bigskip
  \raggedright
  Al leer de una réplica local, cada región puede seguir respondiendo
  aunque se pierda la comunicación con las demás: prioriza la
  disponibilidad (A).

  \vfill
  \footnotesize Fuente: \cite{leenders2018padlock} (documentación técnica
  de Spotify, 2018).
\end{frame}

% ---------------------------------------------------------------------
\begin{frame}{Spotify: decisión documentada sin partición}
  \small
  \renewcommand{\arraystretch}{1.4}
  \begin{tabular}{p{0.30\textwidth} p{0.58\textwidth}}
    \toprule
    \textbf{Aspecto} & \textbf{Documentación de Spotify} \\
    \midrule
    Datos analizados & Personalización, almacenada en Apache Cassandra \\
    Requisito & No requieren transacciones estrictas; pueden
      reconstruirse desde el flujo de escucha \\
    Decisión & Menor consistencia inmediata entre réplicas a cambio de
      menor latencia \\
    Mismo criterio & Padlock tolera unos segundos de retraso en las
      actualizaciones para minimizar la latencia \\
    Clasificación PACELC & Con partición, A; sin partición, L: PA/EL \\
    \bottomrule
  \end{tabular}

  \vfill
  \footnotesize Fuentes: \cite{mishra2015personalization},
  \cite{leenders2018padlock} (documentación técnica de Spotify).
\end{frame}

% ---------------------------------------------------------------------
\begin{frame}{Evidencia técnica: niveles ONE y QUORUM}
  Prueba de carga de Spotify sobre Cassandra con un objetivo de latencia
  inferior a 5~ms en el percentil 95.

  \bigskip
  \centering
  \small
  \renewcommand{\arraystretch}{1.4}
  \begin{tabular}{lll}
    \toprule
    \textbf{Clúster} & \textbf{Nivel ONE} & \textbf{Nivel QUORUM} \\
    \midrule
    3 nodos & No cumple el objetivo & No reportado \\
    6 nodos & Cumple (cerca de 90\,000 operaciones/s) & Cumple \\
    \bottomrule
  \end{tabular}

  \bigskip
  \raggedright
  En esta prueba, ampliar el clúster permitió elevar el nivel de
  consistencia sin superar el límite de latencia.

  \vfill
  \footnotesize Fuente: \cite{cody2015cassandra}.
\end{frame}
